% from ega3-6-fr.tex
\oldpage[III]{161}
$U_\alpha^{(1)}\times_S U_\beta^{(2)}$ et
$\mathscr Q_\bullet
=\mathscr P_\bullet^{(1)}\otimes_S\mathscr P_\bullet^{(2)}$ (considéré comme
complexe simple pour le degré total); en effet, la définition de cet
homomorphisme a en substance été donnée au cours du calcul de la suite $(b')$
dans (6.6.4), pour $q=0$; il suffit simplement (en gardant les notations de
(6.6.4)) de tenir compte de ce qu'il y a d'une part un homomorphisme naturel
du complexe $N^{\bullet k}$ dans le complexe
$C^\bullet(\Phi^{(1)},\Phi^{(2)};\mathscr S_k)$
($0_{\mathrm I}$, 11.8.5), d'autre part un homomorphisme naturel de ce dernier
complexe dans le complexe
$P^\bullet(\Phi^{(1)},\Phi^{(2)};\mathscr S_k)$
($0_{\mathrm I}$, 11.8.6), et enfin un homomorphisme naturel de ce dernier
complexe de cochaînes dans le sous-complexe des cochaînes alternées
($0_{\mathrm I}$, 11.8.7). Par ailleurs, l'homomorphisme de bicomplexes ainsi
défini donne un isomorphisme en homologie, comme on l'a vu en (6.6.4); on a
donc, en composant avec (6.7.6.2), obtenu un isomorphisme
\[
  \mathbf{Tor}_n^S
  (\mathfrak U^{(1)},\mathfrak U^{(2)};
   \mathscr P_\bullet^{(1)},\mathscr P_\bullet^{(2)})
  \xrightarrow{\sim}
  \mathbf H^{-n}
  (\mathfrak U,\mathscr P_\bullet^{(1)}\otimes_S\mathscr P_\bullet^{(2)}).
  \tag{6.7.6.3}\label{III.6.7.6.3.fr}
\]
Il faut ensuite prouver que l'isomorphisme ainsi défini ne dépend pas des
recouvrements ouverts choisis (le second membre de (6.7.6.3) étant
canoniquement isomorphe à
$\mathbf H^{-n}
(X^{(1)}\times_S X^{(2)},
 \mathscr P_\bullet^{(1)}\otimes_S\mathscr P_\bullet^{(2)})$ par (6.2.2));
cela se fait à l'aide de (6.6.7) en remarquant (avec les notations de (6.6.7))
que l'on a un diagramme commutatif à homotopismes près
\[
\xymatrix@C=18pt@R=28pt{
  C^\bullet(\mathfrak U^{(1)},\mathscr P_\bullet^{(1)})
  \otimes_A
  C^\bullet(\mathfrak U^{(2)},\mathscr P_\bullet^{(2)})
  \ar[r]\ar[d]
  & C^\bullet(\mathfrak U,\mathscr Q_\bullet)\ar[d]
  \\
  C^\bullet(\mathfrak V^{(1)},\mathscr P_\bullet^{(1)})
  \otimes_A
  C^\bullet(\mathfrak V^{(2)},\mathscr P_\bullet^{(2)})
  \ar[r]
  & C^\bullet(\mathfrak V,\mathscr Q_\bullet)
}
\]
où les flèches horizontales sont les homomorphismes définis ci-dessus. Enfin,
il faut passer au cas général par recollement, ce qui se fait sans difficulté
comme dans (6.7.1) et (6.7.2); nous laissons les détails au lecteur.
\end{proof}

\begin{proposition}[6.7.7]
Supposons que $\mathscr P_\bullet^{(1)}$ et
$\mathscr P_\bullet^{(2)}$ soient limités inférieurement, et que tous les
Modules $\mathscr H^{-n}(f_1,\mathscr P_\bullet^{(1)})$ ou tous les Modules
$\mathscr H^{-n}(f_2,\mathscr P_\bullet^{(2)})$ soient $S$-plats. On a alors un
isomorphisme canonique de bi-$\partial$-foncteurs ($n$ parcourant $\mathbf Z$)
\[
  \mathfrak{Tor}_n^S
  (f_1,f_2;\mathscr P_\bullet^{(1)},\mathscr P_\bullet^{(2)})
  \xrightarrow{\sim}
  \bigoplus_{q_1+q_2=n}
  \mathscr H^{-q_1}(f_1,\mathscr P_\bullet^{(1)})
  \otimes_S
  \mathscr H^{-q_2}(f_2,\mathscr P_\bullet^{(2)}).
  \tag{6.7.7.1}\label{III.6.7.7.1.fr}
\]
\end{proposition}

\begin{proof}
En effet, vu (6.5.8), la suite spectrale $(a)$ de (6.7.3) est dégénérée, et
la proposition résulte aussitôt de ($0_{\mathrm I}$, 11.1.6), cette suite
étant birégulière (6.7.3).
\end{proof}

\begin{theorem}[6.7.8]
Supposons que : $1^\circ$ les complexes $\mathscr P_\bullet^{(1)}$ et
$\mathscr P_\bullet^{(2)}$ soient limités inférieurement; $2^\circ$ le
complexe $\mathscr P_\bullet^{(1)}$ ou le complexe
$\mathscr P_\bullet^{(2)}$ soit formé de Modules $S$-plats; $3^\circ$ tous les
